\clearpage\begingroup
\emergencystretch=4em\relax\NoSourceHead\SourceLeaf{461}
\makeatletter\tagsleft@true\let\veqno\@@leqno\makeatother
\renewcommand{\thefootnote}{\arabic{footnote})}\setcounter{footnote}{0}
\ArithmeticHead{4}{L}{BOLZA, GAUSS UND DIE VARIATIONSRECHNUNG.}
\begin{center}\bfseries
I. Teil: Die Principia generalia theoriae figurae fluidorum in statu aequilibrii.
\end{center}
\noindent\hspace*{\parindent}In der Arbeit: \textit{Principia generalia theoriae figurae fluidorum in statu aequilibrii}
werden zum ersten Mal in der Geschichte der Variationsrechnung für
ein zweidimensionales Variationsproblem mit variablem Integrationsbereich die
aus dem Verschwinden der ersten Variation folgenden Bedingungen vollständig
aufgestellt, also nicht nur die partielle Differentialgleichung des Problems,
sondern insbesondere auch die auf die Begrenzung bezügliche \textls[100]{Randbedingung}.

Um daher die Bedeutung dieser Arbeit für die Variationsrechnung würdigen
zu können, wird es nötig sein, einen Überblick über die Geschichte der Extrema
von Doppelintegralen vor \textsc{Gauss} vorauszuschicken.

\begin{center}\textls[80]{I. Abschnitt: Überblick über die Geschichte der Extrema}\par
\textls[80]{von Doppelintegralen vor Gauss.}\end{center}
\noindent\hspace*{\parindent}1. Schon in seiner ersten grundlegenden Arbeit vom Jahre 1760/1 hatte
\textsc{Lagrange}\footnote{\textit{Essai d'une nouvelle méthode pour déterminer les maxima et les minima des formules intégrales indéfinies}, Miscell. Taurin., Bd. II, 1760—1761, und Oeuvres de \textsc{Lagrange}, Bd. I, S. 353—356.}
seine neue Variationsmethode auf das Problem der Fläche kleinsten
Inhalts\linebreak \textls[100]{bei gegebener Begrenzung} angewandt und die Aufgabe bis zur
Aufstellung der partiellen Differentialgleichung der Minimalflächen geführt.

In Verallgemeinerung dieser \textsc{Lagrange}schen Untersuchungen hat dann
\textsc{Euler}\footnote{Institutiones calculi integralis, Bd. III, Appendix de calculo variationum (1770), Art. 159—174, Opera, Ser. I, vol. 13 (1914), S. 458—469.}
1770 den ersten Versuch einer allgemeinen Theorie der Extrema
von Doppelintegralen gemacht. Er betrachtet ein Doppelintegral von der Form
\begin{equation*}
J=\iint V(x,y,z,p,q,\ldots)\,dx\,dy\tag{1}
\end{equation*}
wo $z$ eine Funktion der unabhängigen Variabeln $x$, $y$ ist und
\[
p=\frac{\partial z}{\partial x},\quad q=\frac{\partial z}{\partial y},\quad\ldots,
\]
zwar nicht bei fester Begrenzung des gesuchten Flächenstücks, aber doch bei
\textls[100]{festem Integrationsbereich}.
\par\clearpage\SourceLeaf{462}
\ArithmeticHead{5}{R}{GESCHICHTE DER EXTREMA VON DOPPELINTEGRALEN VOR GAUSS.}
\noindent\hspace*{\parindent}Da hier nur die Funktion $z$ variiert zu werden braucht, während sowohl
die Grenzen wie die unabhängigen Variabeln unvariiert bleiben, so bietet die
\textls[100]{Aufstellung} der ersten Variation von $J$ keine Schwierigkeiten. Der in bekannter
Weise erhaltene Ausdruck für $\delta J$ wird dann nach dem Vorgang von
\textsc{Lagrange} mittels der der partiellen Integration zugrunde liegenden Formeln
\begin{equation*}
u\frac{\partial v}{\partial x}=\frac{\partial(uv)}{\partial x}-v\frac{\partial u}{\partial x},\quad
u\frac{\partial v}{\partial y}=\frac{\partial(uv)}{\partial y}-v\frac{\partial u}{\partial y}\tag{2}
\end{equation*}
auf die Form gebracht
\begin{equation*}
\delta J=\iint\Omega\omega\,dx\,dy+\iint\left(\frac{\partial A}{\partial x}+\frac{\partial B}{\partial y}\right)dx\,dy.\tag{3}
\end{equation*}
Darin ist $\omega$ der unendlichkleine Zuwachs von $z$ und
\begin{equation*}
\left\{\begin{aligned}
&\Omega=N-\frac{\partial P}{\partial x}-\frac{\partial Q}{\partial y}\ldots,\\
&A=P\omega+\ldots,\quad B=Q\omega+\ldots,
\end{aligned}\right.\tag{4}
\end{equation*}
wenn die partiellen Ableitungen von $V$ nach $z,p,q,\ldots$ beziehungsweise mit
$N,P,Q,\ldots$ bezeichnet werden.

Durch Ausführung je einer Integration in dem zweiten Doppelintegral
wird dann gefolgert
\begin{equation*}
\iint\left(\frac{\partial A}{\partial x}+\frac{\partial B}{\partial y}\right)dx\,dy=\int A\,dy+\int B\,dx.\tag{5}
\end{equation*}

Aus dem Verschwinden von $\delta J$ wird nun geschlossen, dass das in $\delta J$
übrig bleibende Doppelintegral bei beliebiger Wahl von $\omega$ verschwinden muss,
was zunächst auf die partielle Differentialgleichung des Problems
\begin{equation*}
\Omega=0\tag{6}
\end{equation*}
und weiterhin auf die \textls[100]{Grenzgleichung}
\begin{equation*}
\int A\,dy+\int B\,dx=0\tag{7}
\end{equation*}
führt.

Weitere Schlüsse aus dieser Gleichung (7) zu ziehen ist \textsc{Euler} nicht gelungen;
er sagt darüber\footnote{\textit{Methodus nova et facilis calculum variationum tractandi}, Novi Comment. Petropol. Bd. XVI (1771, 1772, S. 53); abgedruckt in Institutiones calculi integralis, Bd. IV, 3. Aufl. (1845), S. 605.}:
\par\clearpage\SourceLeaf{463}
\ArithmeticHead{6}{L}{BOLZA, GAUSS UND DIE VARIATIONSRECHNUNG.}
\noindent\hspace*{\parindent}»Verum quid haec singula membra (nämlich der Gleichung (7)) proprie
significent et ad quemnam usum adhiberi queant, neutiquam adhuc perspicere
licet, unde hoc argumentum cujus prima fundamenta etiamnunc vix jacta sunt
censenda, omnem geometrarum attentionem atque multo accuratiorem investigationem
postulare videtur, quod negotium vix ante suscipere licet, quam
casus nonnulli particulares omni studio et diligentia fuerint evoluti«.

Und ein anderer guter Kenner der Variationsrechnung, \textsc{Brunacci}, bricht
noch vierzig Jahre später, nachdem er in seinem \textit{Corso di matematica sublime}\footnote{Bd. IV, Florenz 1808, S. 248.}
die \textsc{Euler}schen Untersuchungen reproduziert hat, an derselben Stelle in die
Klage aus:

»Siamo sopra una spiaggia da cui si scopre una mar senza fine, e non ci
è dato per anche d’inoltrarvisi, onde fare delle scoperte«.

In der Tat stehen wir hier vor einer der Hauptschwierigkeiten der Theorie
der Extrema von Doppelintegralen, die fast sechzig Jahre lang den Fortschritt
der Variationsrechnung nach dieser Richtung aufgehalten hat, und die erst
durch \textsc{Gauss} vollständig überwunden worden ist\footnote{Auf eine zweite, mit der partiellen Integration verbundene Schwierigkeit, auf die ebenfalls bereits \textsc{Euler} gestossen ist, und die erst durch \textsc{Poisson} ihre Erledigung gefunden hat, gehen wir hier nicht ein, da sie erst eintritt, wenn unter dem Integral auch die zweiten Ableitungen von $z$ vorkommen, was bei dem \textsc{Gauss}schen Problem nicht der Fall ist. Vgl. hierüber \textsc{Todhunter}, \textit{A history of the progress of the calculus of variations during the nineteenth century}, Cambridge 1861, Art. 27—29.}.

2. Während bei mehrdimensionalen Variationsproblemen mit festem Integrationsbereich
die \textls[100]{Aufstellung} der ersten Variation keinerlei Schwierigkeiten
bietet, diese sich vielmehr erst bei der \textls[100]{Umformung} der ersten
Variation durch partielle Integration einstellen, treten bei Problemen mit
\textls[100]{variablem Integrationsbereich},\linebreak oder, was dasselbe ist, bei variabeln
Grenzen, schon sofort bei der Aufstellung der ersten Variation erhebliche
Schwierigkeiten auf. Zu ihrer Überwindung sind von Anfang an zwei wesentlich
von einander verschiedene Methoden angewandt worden, die sich auch
in ihrer weiteren Entwicklung scharf von einander getrennt gehalten haben:
die Methode der Variation der Grenzen und die Methode der Variation der
unabhängigen Variabeln.

Die Methode der \textls[100]{Variation der Grenzen} ist dadurch charakterisiert,
dass bei ihr von vornherein das mehrfache Integral als Aufeinanderfolge
\par\clearpage\SourceLeaf{464}
\ArithmeticHead{7}{R}{GESCHICHTE DER EXTREMA VON DOPPELINTEGRALEN VOR GAUSS.}
\noindent einfacher Integrale mit fester Integrationsordnung aufgefasst wird, und dass die
unbekannten Funktionen und die Grenzen, nicht aber die unabhängigen
Variabeln variiert werden. Demgegenüber liegt der Methode der \textls[100]{Variation
der unabhängigen Variabeln}\linebreak bewusst oder unbewusst die Auffassung des
mehrfachen Integrals als Grenze einer Summe zugrunde, und es werden die
unbekannten Funktionen und die unabhängigen Variabeln, nicht aber die
Grenzen variiert. Die Methode der Variation der Grenzen ist die nächstliegende,
begrifflich einfachere; sie hat rein analytischen Charakter und ist
schon frühzeitig mit Klarheit und Präzision und fehlerfrei durchgeführt worden.
Die Methode der Variation der unabhängigen Variabeln dagegen hat sich in
näherem Zusammenhang mit der Anschauung gehalten; sie ist jedoch von
ihrem ersten Auftreten bei \textsc{Lagrange} an mit Unklarheiten und Missverständnissen
belastet gewesen, die vielfach zu direkten Fehlern geführt haben.
Diesen Nachteilen stehen aber auch entscheidende Vorteile gegenüber: Die
Methode der Variation der unabhängigen Variabeln ist nicht nur für die Behandlung
geometrischer Aufgaben die geeignetere, sondern man kann geradezu
sagen, dass die Methode der Variation der Grenzen, wenigstens in der Form,
in der sie sich geschichtlich weiterentwickelt hat, für eine befriedigende Lösung
geometrischer Aufgaben untauglich ist (siehe unten Nr. 5).
\par\clearpage\endgroup
